\subsubsection{1D Distributed Area--Flow Models}
\label{subsubsec:1d-models}

The selected 1D reduction is defined in
Section~\ref{subsubsec:cross-sectional-reduction-selected-1d-model}. There the
retained fields $A(z,t)$, $Q(z,t)$, $\bar u(z,t)=Q/A$, and
$p_{\mathrm g}(z,t)$, the wall and coupling closure catalog, the elastic-wall
area--flow balance law, the stenosis reference geometry, the conservative
flux/source notation, and the single-vessel boundary data are recorded in
Definitions~\ref{def:section-averaged-1d-fields}--\ref{def:1d-single-vessel-boundary-data}
\cite{FormaggiaEtAl2003OneDimensionalBloodFlow,KopplHelmig2023DimensionReducedModeling}.
This subsection uses that machinery only to identify the hierarchy role of the
1D tier.

In the hierarchy, the key point is that the 1D state is distributed but reduced.
It retains axial area, flow, mean velocity, and gauge pressure; it does not
retain the full three-dimensional velocity field, pressure field, secondary
flow, separation, recirculation, or wall traction. Reduced shear quantities
therefore inherit the declared velocity-profile, wall-friction, and rheology
closures. They are not traction-based WSS unless a separate observation map
declares how a reduced output is compared with a resolved wall-stress quantity.

The closure contract remains part of the model definition. A 1D case must
declare the momentum-flux/profile convention, friction law, rheology descriptor
$\mathcal R$, wall descriptor $\mathcal W$, boundary data, numerical
realization, and output operators. Stenosis enters through the reference
geometry $R_0(z)$, $A_0(z)$, and $s(z)$, but resolved throat flow does not
appear automatically; variable-radius corrections such as the Canic-style terms
are additional closure choices
\cite{CanicEtAl2024Extended1DStenosis}.
The current implementation records this axis explicitly through the
\texttt{canic-extended-1d} and \texttt{classical-1d-no-slip} forward-model
descriptors, the latter corresponding to the wall no-slip baseline in
Definition~\ref{def:classical-no-slip-1d-baseline}.

The 1D model is the selected reduced forward map because it retains distributed
pressure--flow structure at low computational cost, while its outputs remain
reduced quantities with declared observation maps. In this report it is the
selected single-vessel generator for the numerical study, not a
clinical-validation claim, an independent-solver parity claim, or a
resolved-3D-accuracy claim.
